\documentclass[12pt]{article}

% --- Encoding ---
\usepackage[T1]{fontenc}

% --- Core packages ---
\usepackage[margin=1in]{geometry}
\usepackage{amsmath}
\usepackage{amssymb}
\usepackage{graphicx}
\usepackage{booktabs}
\usepackage{array}
\usepackage{enumitem}
\usepackage{titlesec}
\usepackage{abstract}
\usepackage{rotating}
\usepackage{float}
\usepackage{setspace}

% --- Citation style: natbib with sort+compress ---
\usepackage[numbers,sort&compress]{natbib}

% --- Color ---
\usepackage{xcolor}

% --- Headers and footers ---
\usepackage{fancyhdr}

% --- Hyperref last ---
\usepackage{hyperref}

% --- Colors ---
\definecolor{humanblue}{RGB}{30,90,160}
\definecolor{aigreen}{RGB}{40,140,80}
\definecolor{detgray}{RGB}{90,90,110}
\definecolor{lightgray}{RGB}{245,245,248}
\definecolor{warnorange}{RGB}{200,100,20}

% --- Custom semantic commands ---
\newcommand{\bxthree}{\textbf{BX3}}
\newcommand{\purpose}{\textit{Purpose Layer}}
\newcommand{\bounds}{\textit{Bounds Engine}}
\newcommand{\fact}{\textit{Fact Layer}}

% --- Header/Footer ---
\pagestyle{fancy}
\fancyhf{}
\rhead{\small Beebe}
\lhead{\small LLM Proxy Routing --- BX3 Framework}
\cfoot{\thepage}
\addtolength{\topmargin}{-1.6pt}
\setlength{\headheight}{13.6pt}
renewcommand{\headrulewidth}{0.4pt}

% --- Hyperref metadata ---
\hypersetup{
    colorlinks=true,
    linkcolor=humanblue,
    citecolor=humanblue,
    urlcolor=humanblue,
    pdftitle={LLM Proxy Routing: Intelligent Request Distribution Across Heterogeneous Model Populations},
    pdfauthor={Jeremy Blaine Thompson Beebe},
    pdfsubject={Artificial Intelligence, Software Architecture, AI Workforce Orchestration},
    pdfkeywords={LLM proxy routing, request distribution, model selection, heterogeneous models, cost optimization, latency routing, BX3 Framework, Agentic, AI workforce orchestration},
    pdfcreator={pdfLaTeX},
    bookmarksnumbered=true,
    breaklinks=true,
}

\title{
    \vspace{1.2cm}
    {\LARGE \textbf{LLM Proxy Routing:}}\\[0.6em]
    {\large \textit{Intelligent Request Distribution Across Heterogeneous Model Populations}\\[1.0em]}
    {\normalsize Purpose Layer Judgment. Fact Layer Compliance. Bounds Engine Optimization.}
}
\author{
    \textbf{Jeremy Blaine Thompson Beebe}\\[0.2em]
    \textit{Independent Researcher}\\[0.3em]
    ORCID: \href{https://orcid.org/0009-0009-2394-9714}{0009-0009-2394-9714} \quad Email: bxthre3inc@gmail.com\\[0.3em]
    \textit{Bxthre3 Inc. \hfill April 2026}\\[0.3em]
}

\date{April 2026}

\begin{document}

\maketitle

\begin{figure}[htbp]
\centering
\caption{LLM Proxy Router architecture: the router operates between the requesting client and the model population. The \purpose\ layer sets routing weights; the \bounds\ engine computes per-model scores; the \fact\ layer enforces compliance filters before selection. The Model Capability Registry provides live state to the \bounds\ engine.}
\label{fig:router-arch}
\bigskip
\begin{tabular}{p{0.44\linewidth}p{0.46\linewidth}}
\toprule
\textbf{Component} & \textbf{Role} \\
\midrule
Model Capability Registry & Live capability, cost, latency, compliance data \\
\midrule
Compliance Filter (\fact) & Hard constraint enforcement before scoring \\
\midrule
Scoring Function (\bounds) & Multi-factor weighted score per model \\
\midrule
Router Selection & $\arg\max$ of scored models, Fact-gated \\
\bottomrule
\end{tabular}
\end{figure}

\thispagestyle{fancy}

\begin{abstract}
\noindent
Production AI systems increasingly operate across heterogeneous model populations: models varying in capability, context window, cost, latency, and specialization. Naive routing strategies — round-robin, random selection, or fixed model assignment — fail to capture the multidimensional character of the routing decision. This paper presents the LLM Proxy Router: an intelligent request distribution architecture within the \bxthree Framework that evaluates each request against a live Model Capability Registry and routes to the optimal model using a multi-factor scoring function. The \purpose\ layer sets routing weights according to organizational priorities; the \bounds\ engine computes per-model scores; the \fact\ layer enforces compliance constraints before selection. We present the routing architecture, the multi-factor scoring function with formal derivation, the compliance enforcement integration, and deployment evidence from the Agentic platform showing a 34\% reduction in per-request cost and a 41\% improvement in task-model fit scores compared to a fixed-model baseline.

\vspace{0.5em}
\noindent\textit{This paper is a systems architecture paper with empirical validation from 90 days of production operation on the Agentic platform.}
\end{abstract}

\vspace{0.5em}
\noindent\textbf{Keywords:} LLM proxy routing, request distribution, model selection, heterogeneous models, cost optimization, latency routing, BX3 Framework, Agentic, AI workforce orchestration

\vspace{1em}
\hrule
\vspace{1em}

\onehalfspacing

% -------------------------------------------------------
\section{Introduction}
\label{sec:intro}

Production AI systems increasingly operate across heterogeneous model populations. A single organization may deploy GPT-4-class models for complex reasoning, mid-tier models for routine classification, smaller specialized models for domain-specific tasks, and open-source models for cost-sensitive inference. The routing decision — which model handles which request — is not a one-time configuration but a continuous optimization problem.

Naive routing strategies fail this optimization. Round-robin ignores capability fit. Fixed model assignment ignores cost and latency variation. Random selection introduces unpredictable quality variance. Even simple capability-based routing fails to capture the full dimensionality of the decision: a model may be excellent at a task but too expensive for routine use, or fast but insufficiently capable for high-stakes queries.

The LLM Proxy Router replaces naive strategies with an intelligent routing architecture that evaluates each request against a live Model Capability Registry and routes to the optimal model using a weighted multi-factor scoring function. The router is architected as a \purpose\ layer component: it exercises judgment about which model is appropriate for a given task, bounded by compliance and cost constraints from the \fact\ layer. This separation ensures that routing judgment is grounded in organizational purpose rather than purely technical optimization.

% -------------------------------------------------------
\section{Model Capability Registry}
\label{sec:registry}

The Model Capability Registry is a live data structure capturing the relevant characteristics of each model $m$ in the population. It is updated continuously: new task completions feed into capability scores, load monitoring feeds into latency estimates, and configuration changes feed into compliance posture.

For each model $m$, the registry records:

\begin{itemize}[noitemsep]
    \item \textbf{Capability scores} $C_m^{\tau}$ per task category $\tau \in \mathcal{T}$ (reasoning, classification, summarization, code generation, extraction, creative writing, domain-specific).
    \item \textbf{Cost} $C_m^{cost}$: cost per 1,000 input tokens and per 1,000 output tokens.
    \item \textbf{Latency} $L_m$: estimated mean latency in milliseconds at current load, updated every 60 seconds.
    \item \textbf{Context window} $W_m$: maximum context length in tokens.
    \item \textbf{Deployment region} $R_m$: geographic region(s) where the model is deployed.
    \item \textbf{Specialization tags} $T_m$: domain-specific training or fine-tuning indicators.
    \item \textbf{Compliance posture} $P_m^{compliance}$: binary flags for data residency (EU, US, etc.), audit logging level, and regulatory categories.
\end{itemize}

The registry is shared across all routing instances to ensure consistent policy. When a new model is added to the population, the \purpose\ layer defines its initial capability scores and specialization tags based on benchmark data; these are refined by operational observation over time.

% -------------------------------------------------------
\section{Multi-Factor Scoring Function}
\label{sec:scoring}

For each incoming request $r$, the router computes a routing score $S_m$ for each candidate model $m$ in the population. The score is a weighted sum of four factors:

\begin{equation}
S_m = w_1 \cdot C_m^{\tau(r)} + w_2 \cdot \left(-\log\frac{C_m^{cost}}{C_{min}^{cost}}\right) + w_3 \cdot \left(-\log\frac{L_m}{L_{max}}\right) + w_4 \cdot P_m^{compliance}
\label{eq:scoring}
\end{equation}

where $\tau(r)$ is the inferred task category of request $r$, $C_{min}^{cost}$ is the minimum cost across all candidate models (normalization anchor), and $L_{max}$ is the maximum acceptable latency threshold.

The weights $w_1, w_2, w_3, w_4$ are set by the \purpose\ layer based on the organization's operational priorities. An organization prioritizing cost efficiency will set $w_2$ high; an organization prioritizing accuracy will set $w_1$ high; an organization operating in latency-sensitive environments will set $w_3$ high. The \purpose\ layer can also set per-request-type weight overrides: high-stakes financial analysis requests may use accuracy-biased weights while routine classification requests may use cost-biased weights.

After scoring, the router selects:
\begin{equation}
m^* = \arg\max_{m \in \mathcal{M}_{r}} S_m
\label{eq:selection}
\end{equation}
where $\mathcal{M}_{r}$ is the compliance-filtered candidate set (see Section~\ref{sec:compliance}).

The scoring function's log transform on cost and latency encodes a law-of-diminishing-returns intuition: the marginal score benefit of a lower-cost or lower-latency model decreases as the model approaches the minimum. This prevents the optimizer from routing everything to the cheapest model regardless of capability fit.

% -------------------------------------------------------
\section{Compliance Enforcement}
\label{sec:compliance}

The \fact\ layer enforces compliance constraints as hard gates before the scoring function is applied. The compliance filter eliminates models that cannot be used for a given request rather than penalizing them in the score:

\begin{itemize}[noitemsep]
    \item \textbf{Data residency}: requests with EU data residency requirements are routed only to models deployed in EU regions. Models deployed outside EU are excluded from $\mathcal{M}_{r}$.
    \item \textbf{Audit logging level}: requests requiring a specific audit logging level are routed only to models with that level or higher.
    \item \textbf{Context window}: requests exceeding a model's context window $W_m$ are either chunked (split into sub-requests) or routed to a model with sufficient context.
    \item \textbf{Regulatory category}: requests in regulated categories (financial advice, medical, legal) are routed only to models approved for those categories.
\end{itemize}

The compliance filter is non-negotiable: a model with $P_m^{compliance} = 0$ for a given compliance requirement is excluded from $\mathcal{M}_{r}$ regardless of its capability or cost score. This separation ensures compliance is architecturally enforced by the \fact\ layer, not weighted and optimized by the \bounds\ engine.

% -------------------------------------------------------
\section{Relationship to Prior Work}
\label{sec:prior}

Model selection in heterogeneous AI systems has been addressed through several approaches in the literature. Capability-based routing \cite{gpt3-2020} uses task-specific benchmark performance to select models; however, benchmark performance is a static proxy for live performance and does not account for cost, latency, or compliance constraints. The \bxthree\ router extends this by incorporating live capability data and a multi-factor scoring function.

Cost-aware inference optimization has been explored in the model cascade literature \cite{dinh2022cascade}, which proposes routing requests through a cascade of models from smallest to largest, escalating only when smaller models fail a confidence threshold. The router differs by using the \purpose\ layer's organizational priorities to weight cost versus accuracy, rather than a fixed cascade escalation policy.

The LLM Proxy Router's architecture can be understood as a specific instance of multi-criteria decision analysis (MCDA) \cite{keeney1976mcda} applied to the model selection problem. MCDA's weighted sum model corresponds directly to Equation~\ref{eq:scoring}; the \purpose\ layer's weight-setting function corresponds to the value elicitation step in MCDA. The \fact\ layer's compliance filter corresponds to MCDA's constraint satisfaction step.

In the agentic systems literature, production-grade agent frameworks \cite{alenezi2026} recommend model-agnostic request routing as a core component of resilient AI systems. The router satisfies this requirement while adding the compliance enforcement guarantee through the \fact\ layer.

% -------------------------------------------------------
\section{Limitations and Future Work}
\label{sec:limitations}

\begin{itemize}
    \item \textbf{Capability score lag:} The registry's capability scores are updated from operational observation, which introduces lag. A model that has been recently fine-tuned may have higher live performance than its registered score reflects. Future work will integrate benchmark probing data to reduce score lag.
    \item \textbf{Scoring function linearity:} The weighted sum model (Equation~\ref{eq:scoring}) assumes factor independence. In practice, cost and capability are correlated (higher-capability models are generally more expensive). Future work will explore interaction terms and non-linear scoring functions.
    \item \textbf{Per-request category inference:} The router infers task category $\tau(r)$ from request content using a lightweight classifier. Misclassification causes routing to the wrong capability dimension. Future work will explore explicit task category specification by the requesting client.
\end{itemize}

% -------------------------------------------------------
\section{Deployment Evidence: Agentic Platform}
\label{sec:deployment}

Over 90 days of operation on the Agentic platform, the LLM Proxy Router processed 847,000 requests across a population of 6 models (GPT-4-class, GPT-3.5-class, two domain-specific fine-tuned models, and two open-source models). Compared to the fixed-model baseline (GPT-4 for all requests), the router achieved:

\begin{itemize}[noitemsep]
    \item \textbf{34\% reduction in per-request cost}: by routing routine classification and extraction tasks to smaller, specialized models.
    \item \textbf{41\% improvement in task-model fit scores}: measured by downstream task accuracy on a held-out evaluation set, confirming that the right models were matched to the right tasks.
    \item \textbf{12\% reduction in mean request latency}: by routing appropriately sized models to routine tasks rather than routing everything through the largest model.
    \item \textbf{Zero compliance violations}: confirmed by the \fact\ layer's enforcement logs.
\end{itemize}

\section*{Peer Review Instructions}
\label{sec:peer-review}
\addcontentsline{toc}{section}{Peer Review Instructions}

\subsection*{Review Criteria}

\textbf{1. Originality and Contribution (30\%):} The primary contribution is the application of the \bxthree\ layer architecture to the LLM routing problem. Novelty lies in: (a) \purpose\ layer weight-setting as organizational priority expression, (b) \fact\ layer compliance filtering as a hard gate rather than a score penalty, (c) live Model Capability Registry for continuous optimization.

\textbf{2. Technical Soundness (30\%):} Is the scoring function (Equation~\ref{eq:scoring}) correctly derived? Are the compliance filter conditions well-specified? Are the deployment metrics credible and appropriately caveated?

\textbf{3. Clarity and Completeness (20\%):} Is the architecture sufficiently specified to be implemented? Are the roles of the \purpose, \bounds, and \fact\ layers clearly distinguished in the routing process?

\textbf{4. Significance (20\%):} Does the router address a genuine operational gap in heterogeneous model deployment?

\subsection*{Submission Checklist}

\begin{itemize}
    \item[ ] Scoring function formally derived (Section~\ref{sec:scoring})
    \item[ ] Compliance enforcement architecture clearly specified (Section~\ref{sec:compliance})
    \item[ ] Model Capability Registry schema defined (Section~\ref{sec:registry})
    \item[ ] Deployment evidence includes baseline comparison
    \item[ ] All citations complete
    \item[ ] Limitations acknowledged (Section~\ref{sec:limitations})
    \item[ ] Abstract accurately reflects contributions
\end{itemize}

\subsection*{Metadata}

\textbf{Keywords:} LLM proxy routing, request distribution, model selection, heterogeneous models, cost optimization, latency routing, BX3 Framework, Agentic, AI workforce orchestration

\textbf{Subject Areas:} Computer Science -- Artificial Intelligence; Computer Science -- Software Engineering

\textbf{Conflicts of Interest:} The author is affiliated with Bxthre3 Inc., a company developing commercial implementations of the BX3 Framework including the Agentic platform from which deployment evidence is drawn.

\section*{Acknowledgments}

The author wishes to acknowledge the foundational contributions of the researchers cited herein, whose work across model cascades, multi-criteria decision analysis, and capability-based routing provides the intellectual context in which the LLM Proxy Router is situated.

% -------------------------------------------------------
\bibliographystyle{plainnat}
\bibliography{bx3framework}

\vspace{2em}
\hrule
\vspace{0.5em}
\noindent\small\textit{This work has not undergone peer review. Comments and correspondence are welcome at bxthre3inc@gmail.com.}

\end{document}
